%%%%%%%%%%%%%%%%%%%%%%%%%%%%%%%%%%%%%%%%%%%%%%%%%%%%%%%%%%%%%%%%%%%%%%%%%%%%%%%%
%                                                                              %
%                            BlackArch Linux Guide                             %
%                                 Spanish version                              %
%                                                   by Ubidragon               %
%                                                                              %
%%%%%%%%%%%%%%%%%%%%%%%%%%%%%%%%%%%%%%%%%%%%%%%%%%%%%%%%%%%%%%%%%%%%%%%%%%%%%%%%
%
%
%
%
%Traductors:
%Ubaldo Hidalgo Arriaga <ubidragon@protonmail.com>, 2016
%
%
%
%

\documentclass[a4paper, oneside, 11pt]{book}

%%% INCLUDES %%%
\renewcommand{\familydefault}{\sfdefault}

\usepackage{array}
\usepackage{color}
\usepackage{enumerate}
\usepackage{fancyhdr}
\usepackage{fancyvrb}
\usepackage{geometry}
\usepackage{graphicx}
\usepackage{html}
\usepackage{hyperref}
\usepackage{ifpdf}
\usepackage{listings}
\usepackage{pstricks}
\usepackage{supertabular}
\usepackage{tocloft}
\usepackage[utf8]{inputenc}

%%% LAYOUT %%%
\setlength{\parindent}{0em}
\setlength{\parskip}{1.5ex plus0.5ex minus0.5ex}
\geometry{left=2.5cm, textwidth=16cm, top=3cm, textheight=25cm, bottom=3cm}
\widowpenalty=2000
\clubpenalty=1000
\frenchspacing
\sloppy
\pagecolor[HTML]{FFFFFF}
\color[HTML]{333333}
\setcounter{tocdepth}{10}
\setcounter{secnumdepth}{10}

\hypersetup{
  pdftitle={BlackArch Linux, The BlackArch Linux Guide},
  pdfsubject={BlackArch Linux, The BlackArch Linux Guide},
  pdfauthor={BlackArch Linux, BlackArch Linux},
  pdfkeywords={BlackArch Linux, Penetration Testing, Security, Hacking, Linux},
  pdfcenterwindow=true,
  colorlinks=true,
  breaklinks=true,
  linkcolor=blue,
  menucolor=blue,
  urlcolor=blue
}

% syntax highlighting
% all options should be set here document wide
\lstset{
backgroundcolor=\color[HTML]{EEEEEE},
frame=single,
basicstyle=\footnotesize\ttfamily,
float,
deletekeywords={return},
otherkeywords={mkdir, curl, sudo, sha1sum, grep, cut, sort, wget, makepkg,
pacman, blackman, chmod},
keywordstyle=\color{orange},
commentstyle=\color{blue},
stringstyle=\color{red},
language=bash,
showspaces=false,
showtabs=false,
tabsize=2
}

%%% HEADER / FOOTER %%%
\setlength{\headheight}{33pt}
\setlength{\headsep}{33pt}
\lhead{{\includegraphics[width=1cm,height=1cm]{images/logo.png}}}
\rhead{The BlackArch Linux Guide}

%%% CUSTOM MACROS %%%
% for html links
\ifpdf\else
\def\href#1#2{\htmladdnormallink{#2}{#1}}
\fi

%------------------%
%  TITLE PAGE      %
%------------------%
\begin{document}
\pagestyle{empty}
\begin{center}
\begin{figure}[htbp]
\centering
\vspace{0.5cm}
\includegraphics[width=8cm]{images/logo.png}
\label{fig:logo}
\end{figure}
\vspace{0.5cm}
\Huge{\textbf{The BlackArch Linux Guide}}\\
\vspace{0.5cm}
\Large{\textbf{Spanish Version}}\\
\vspace{1cm}
\Large{\color{blue}https://www.blackarch.org/}\\
\vspace{0.5cm}
\end{center}
\newpage
\tableofcontents
\newpage
\pagestyle{fancy}

%------------------%
%  Capitulo 1      %
%------------------%

\chapter{Introducci\'on}

\section{Resumen}
The BlackArch Linux Guide se divide en varias partes:
\begin{itemize}
\item Introducci\'on - Proporciona una visi\'on general, introducci\'on, y una informaci\'on adicional sobre el proyecto
\item Gu\'ia del Usuario - Todo lo que un usuario necesita saber para poder utilizar con eficacia BlackArch
\item Gu\'ia del Desarrollador - ¿C\'omo empezar a desarrollar y contribuir para BlackArch?
\item Gu\'ia de Herramientas - Detalles en profundidad de las herramientas con ejemplos de uso (WIP)
\end{itemize}

\section{¿Qu\'e es BlackArch Linux?}
BlackArch es una completa distribuci\'on de Linux para penetration testers e investigadores de seguridad.
Es una distribuci\'on derivada de \href{https://www.archlinux.org/}{ArchLinux} por lo que los
usuarios podr\'an instalar componentes individuales o grupos de componentes directamente sobre
esta distribuci\'on.

La herramienta se distribuye como un Arch Linux
\href{https://wiki.archlinux.org/index.php/Unofficial\_User\_Repositories}
{repositorio no oficial} por lo que podemos instalar BlackArch sobre una instalaci\'on
existente de Arch Linux. Los paquetese se pueden instalar individualmente o por categorias.


El repositorio en constante expansi\'on incluye \href{https://www.blackarch.org/tools.html}{2600} herramientas.
Todas las herramientas son probadas antes de ser agregadas al c\'odigo base para mantener la calidad del repositorio.
% should quickly describe the testing methods/code review procedures etc

\section{Historia de BlackArch Linux}
Pr\'oximamente...

\section{Plataformas Soportadas}
Pr\'oximamente...

\section{Participaci\'on}
Usted puede ponerse en contacto con el equipo de BlackArch por los siguientes canales:

Website: \url{https://www.blackarch.org/}

Email: \href{mailto:team@blackarch.org}{team@blackarch.org}

IRC: \url{irc://irc.freenode.net/blackarch}

Twitter: \url{https://twitter.com/blackarchlinux}

Github: \url{https://github.com/Blackarch/}

Discord: \url{https://discord.com/invite/xMHt8dW}

%------------------%
%  Capitulo 2      %
%------------------%


\chapter{Gu\'ia del Usuario}

\section{Instalaci\'on}
Las siguientes secciones le mostrar\'an c\'omo configurar el repositorio de BlackArch e
instalar paquetes. BlackArch admite instalar desde el repositorio usando paquetes
binarios, as\'i como compilar e instalar de las fuentes.

BlackArch es compatible con una instalaci\'on normal de Arch. Su funcionamiento es
como un repositorio de usuario no oficial. En cambio si usted desea una ISO, mire 
en la secci\'on \href{https://www.blackarch.org/downloads.html#iso}{Live ISO}.

\subsection{Instalaci\'on sobre ArchLinux}
Ejecute \href{https://blackarch.org/strap.sh}{strap.sh} como administrador y sigue
las instrucciones. Mire el ejemplo siguiente.
\begin{lstlisting}
   curl -O https://blackarch.org/strap.sh
   sha1sum strap.sh # should match: d062038042c5f141755ea39dbd615e6ff9e23121
   sudo chmod +x strap.sh
   sudo ./strap.sh
\end{lstlisting}
Ahora descarga una copia nueva de la lista principal de paquetes y paquetes
de sincronizaci\'on:
\begin{lstlisting}
  sudo pacman -Syyu
\end{lstlisting}


\subsection{Instalando Paquetes}
Ahora puede instalar las herramientas del repositorio de BlackArch.
\begin{enumerate}
\item Para ver una lista con todas las herramientas disponibles, ejecute 
\begin{lstlisting}
  pacman -Sgg | grep blackarch | cut -d' ' -f2 | sort -u
\end{lstlisting}

\item Para instalar todas las herramientas, ejecute
\begin{lstlisting}
  pacman -S blackarch
\end{lstlisting}

\item Para instalar una categor\'ia, ejecute
\begin{lstlisting}
  pacman -S blackarch-<category>
\end{lstlisting}

\item Para ver las categor\'ias de BlackArch, ejecute
\begin{lstlisting}
  pacman -Sg | grep blackarch
\end{lstlisting}

\end{enumerate}

\subsection{Instalando paquetes desde fuentes}
Como parte de un m\'etodo alternativo de instalaci\'on, usted puede construir los paquetes
BlackArchpackages desde las fuentes. Usted deber\'a buscar los PKGBUILDs en
\href{https://github.com/BlackArch/blackarch/tree/master/packages}{github}. Para construir
el repositorio entero, usted debr\'a usar la herramienta
\href{https://github.com/BlackArch/blackman}{Blackman}.
\begin{itemize}
\item En primer lugar, usted tiene que instalar Blackman. Si el repositorio de paquetes de BlackArch
est\'a configurado en su equipo, puede instalar Blackman:
\begin{lstlisting}
  pacman -S blackman
\end{lstlisting}

\item Usted puede compilar e instalar desde las fuentes:
\begin{lstlisting}
  mkdir blackman
  cd blackman
  wget https://raw.github.com/BlackArch/blackarch/master/packages/blackman/PKGBUILD
  # Aseg\'urese que PKGBUILD no ha sido manipulado maliciosamente.
  makepkg -s
\end{lstlisting}

\item Tambi\'en puede instalar Blackman desde AUR:
\begin{lstlisting}
  <cualquier ayudante de AUR que use> -S blackman
\end{lstlisting}

\end{itemize}

\subsection{Uso b\'asico de Blackman} Blackman es muy f\'acil de usar, aunque los 
comandos son diferentes de lo que normalmente podr\'ia esperar de algo parecido
a pacman. El uso b\'asico ser\'a delineado a continuaci\'on.
\begin{itemize}
\item Descargue, compile e instale paquetes:
\begin{lstlisting}
  sudo blackman -i package
\end{lstlisting}

\item Descargue, compile e instale alguna categor\'ia:
\begin{lstlisting}
  sudo blackman -g group
\end{lstlisting}

\item Descargue, compile e instale todas las herramientas de BlackArch:
\begin{lstlisting}
  sudo blackman -a
\end{lstlisting}

\item Para ver una lista de las categor\'ias de blackarch:
\begin{lstlisting}
  blackman -l
\end{lstlisting}

\item Para una lista de herramientas por categor\'ias:
\begin{lstlisting}
  blackman -p category
\end{lstlisting}

\end{itemize}

\subsection{Instalando desde live-, netinstall- ISO o ArchLinux}
Usted puede instalar BlackArch Linux desde uno de nuestros live- o netinstall-ISOs.\\Mire
\url{https://www.blackarch.org/download.html#iso}. Los pasos siguientes son necesarios
tras el arranque de la ISO.
\begin{itemize}
\item Instale el paquete blackarch-installer:
\begin{lstlisting}
  sudo pacman -S blackarch-installer
\end{lstlisting}

\item Ejecute
\begin{lstlisting}
  sudo blackarch-install
\end{lstlisting}

\end{itemize}

%------------------%
%  Capitulo 3      %
%------------------%

\chapter{Gu\'ia del Desarrollador}

\section{Sistema de compilaci\'on y repositorios de Arch}

Los archivos PKGBUILD son scripts de construcci\'on. Cada uno makepkg(1) dice c\'omo crear un
paquete. Los archivos PKGBUILD est\'an escritos en Bash.

Para obtener m\'as informaci\'on, lea (u hoje\'e) lo siguiente:
\begin{itemize}
\item \href{https://wiki.archlinux.org/index.php/Creating_Packages}{Arch Wiki: Creando Paquetes}
\item \href{https://wiki.archlinux.org/index.php/Makepkg}{Arch Wiki: makepkg}
\item \href{https://wiki.archlinux.org/index.php/PKGBUILD}{Arch Wiki: PKGBUILD}
\item \href{https://wiki.archlinux.org/index.php/Arch_Packaging_Standards}{Arch Wiki: Arch Packaging Standards}
\end{itemize}

\section{Blackarch PKGBUILD standards}
En favor de la simplicidad, nuestros PKGBUILDs son similares a los de los AUR,
con algunas pequeñas diferencias que se detallan a continuaci\'on. Cada paquete debe
pertenecer a Blackarch como m\'inimo, tambi\'en habr\'a mucho cruce con
m\'ultiples paquetes pertenecientes a m\'ultiples grupos.

\subsection{Grupos}
Para permitir a los usuarios instalar una gama espec\'ifica de paquetes de forma r\'apida y sencilla,
se han separado en grupos. Los grupos permiten a los usuarios simplemente
ir "pacman -S <nombre de grupo>" para sacar un mont\'on de paquetes.

\subsubsection{blackarch}
El grupo de blackarch es el grupo base donde todos los paquetes al que deben pertenecer tambi\'en. Esto permite que
los usuarios puedan instalar cada paquete con facilidad.

Qu\'e deber\'ia estar aqu\'i: Todo.

\subsubsection{blackarch-anti-forensic}
Son paquetes que se utilizan para contrarrestar las actividades forenses,
incluyendo cifrado, esteganograf\'ia y todo lo que modifica los atributos de archivos/archivos.
Todo esto incluye herramientas para trabajar con cualquier cosa en general que realice cambios a un sistema con
el fin de ocultar informaci\'on.

Ejemplos: luks, TrueCrypt, Timestomp, dd, ropeadope, secure-delete

\subsubsection{blackarch-automation}
Paquetes que se utilizan para la automatizaci\'on de herramientas o flujos de trabajo.

Ejemplos: blueranger, tiger, wiffy

\subsubsection{blackarch-backdoor}
Paquetes que explotan o abren puertas traseras en sistemas ya vulnerables.

Ejemplos: backdoor-factory, rrs, weevely

\subsubsection{blackarch-binary}
Paquetes que operan con archivos binarios de alguna forma.

Ejemplos: binwally, packerid

\subsubsection{blackarch-bluetooth}
Paquetes que explotan cualquier cosa relacionada con el est\'andar Bluetooth (802.15.1).

Ejemplos: ubertooth, tbear, redfang

\subsubsection{blackarch-code-audit}
Paquetes que auditan el c\'odigo fuente existente para el an\'alisis de vulnerabilidades.

Ejemplos: flawfinder, pscan

\subsubsection{blackarch-cracker}
Paquetes usados para crackear funciones criptogr\'aficas, i.e.: hashes.

Ejemplos: hashcat, john, crunch

\subsubsection{blackarch-crypto}
Paquetes que trabajan con la criptograf\'ia, a excepci\'on del cracking.

Ejemplos: ciphertest, xortool, sbd

\subsubsection{blackarch-database}
Paquetes que impliquen explotaci\'on de base de datos a cualquier nivel.

Ejemplos: metacoretex, blindsql

\subsubsection{blackarch-debugger}
Paquetes que permiten al usuario ver lo que un programa en particular est\'a "haciendo" en tiempo real.

Ejemplos: radare2, shellnoob

\subsubsection{blackarch-decompiler}
Paquetes que intentan revertir un programa compilado en c\'odigo fuente.

Ejemplos: flasm, jd-gui

\subsubsection{blackarch-defensive}
Paquetes que se utilizan para proteger a un usuario de malware y ataques de otros usuarios.

Ejemplos: arpon, chkrootkit, sniffjoke

\subsubsection{blackarch-disassembler}
Esto es similar a blackarch-decompiler, y probablemente habr\'a mucho
de programas que se encuentran en ambos, sin embargo, estos paquetes producen salida en ensamblador en lugar del c\'odigo fuente en bruto.

Ejemplos: inguma, radare2

\subsubsection{blackarch-dos}
Paquetes que utilizan ataques de denegaci\'on de servicio (DoS).

Ejemplos: 42zip, nkiller2

\subsubsection{blackarch-drone}
Los paquetes que se utilizan para la gesti\'on de la ingenier\'ia f\'isica de drones.

Ejemplos: meshdeck, skyjack

\subsubsection{blackarch-exploitation}
Paquetes que aprovechan las ventajas de los exploits en otros programas o servicios.

Ejemplos: armitage, metasploit, zarp

\subsubsection{blackarch-fingerprint}
Paquetes que explotan los equipos biom\'etricos de huellas dactilares.

Ejemplos: dns-map, p0f, httprint

\subsubsection{blackarch-firmware}
Paquetes que aprovechan las vulnerabilidades del firmware

Ejemplos: None yet, amend asap.

\subsubsection{blackarch-forensic}
Paquetes que se utilizan para encontrar datos en discos f\'isicos o memoria incrustada.

Ejemplos: aesfix, nfex, wyd

\subsubsection{blackarch-fuzzer}
Paquetes que utilizan el principio de prueba de la polic\'ia, es decir
"lanzar" entradas aleatorias al sujeto para ver qu\'e pasa.

Ejemplos: msf, mdk3, wfuzz

\subsubsection{blackarch-hardware}
Paquetes que explotan o gestionan cualquier cosa que tenga que ver con
hardware f\'isico.

Ejemplos: arduino, smali

\subsubsection{blackarch-honeypot}
Los paquetes que act\'uan como "honeypots", es decir, programas que parecen
ser servicios vulnerables utilizados para atraer a los hackers a una trampa.

Ejemplos: artillery, bluepot, wifi-honey

\subsubsection{blackarch-keylogger}
Paquetes que registran y conservan las pulsaciones de teclas en otro sistema.

Ejemplos: None yet, amend asap.

\subsubsection{blackarch-malware}
Paquetes que cuentan como cualquier tipo de software malicioso o
detecci\'on de malware.

Ejemplos: malwaredetect, peepdf, yara

\subsubsection{blackarch-misc}
Paquetes que no encajan particularmente en ninguna categor\'ia.

Ejemplos: oh-my-zsh-git, winexe, stompy

\subsubsection{blackarch-mobile}
Paquetes que manipulan plataformas m\'oviles.

Ejemplos: android-sdk-platform-tools, android-udev-rules

\subsubsection{blackarch-networking}
Paquete que involucra redes IP.

Ejemplos: Anything pretty much

\subsubsection{blackarch-nfc}
Paquetes que utilizan nfc (comunicaciones de campo cercano).

Ejemplos: nfcutils

\subsubsection{blackarch-packer}
Paquetes que operan o invaden a los empaquetadores.

\textit{son programas que incrustan malware dentro de otros ejecutables.}

Ejemplos: packerid

\subsubsection{blackarch-proxy}
Paquetes que act\'uan como proxy, es decir, redirigiendo el tr\'afico
a trav\'es de otro nodo en Internet.

Ejemplos: burpsuite, ratproxy, sslnuke

\subsubsection{blackarch-recon}
Paquetes que buscan activamente vulnerabilidades a
lo salvaje. M\'as bien un grupo de cobertura para paquetes similares.

Ejemplos: canri, dnsrecon, netmask

\subsubsection{blackarch-reversing}
Este es un grupo de cobertura para cualquier descompilador,
desensamblador o cualquier otro programa similar.

Ejemplos: capstone, radare2, zerowine

\subsubsection{blackarch-scanner}
Paquetes que analizan los sistemas seleccionados en busca de vulnerabilidades.

Ejemplos: scanssh, tiger, zmap

\subsubsection{blackarch-sniffer}
Paquetes que involucran el an\'alisis del tr\'afico de la red.

Ejemplos: hexinject, pytactle, xspy

\subsubsection{blackarch-social}
Paquetes que atacan principalmente los sitios de redes sociales.

Ejemplos: jigsaw, websploit

\subsubsection{blackarch-spoof}
Paquetes que intentan suplantar al atacante, tales como
el atacante no parecer como atacante a la v\'ictima.

Ejemplos: arpoison, lans, netcommander

\subsubsection{blackarch-threat-model}
Paquetes que se utilizar\'ian para informar/grabar el
modelo de amenaza descrito en un escenario particular.

Ejemplos: magictree

\subsubsection{blackarch-tunnel}
Los paquetes que se utilizan para tunelizar el tr\'afico de red en una red determinada
red.
Ejemplos: ctunnel, iodine, ptunnel

\subsubsection{blackarch-unpacker}
Los paquetes que se utilizan para extraer malware preempaquetado de un archivo
ejecutable.

Ejemplos: js-beautify

\subsubsection{blackarch-voip}
Paquetes que operan con programas y protocolos voip.

Ejemplos: iaxflood, rtp-flood, teardown

\subsubsection{blackarch-webapp}
Paquetes que funcionan con aplicaciones orientadas a Internet.

Ejemplos: metoscan, whatweb, zaproxy

\subsubsection{blackarch-windows}
Este grupo es para cualquier paquete nativo de Windows que se ejecute v\'ia wine.

Ejemplos: 3proxy-win32, pwdump, winexe

\subsubsection{blackarch-wireless}
Paquetes que operan en redes inal\'ambricas en cualquier nivel.

Ejemplos: airpwn, mdk3, wiffy

\section{Estructura del repositorio}
Usted puede encontrar el repositorio git principal de  BlackArch aqu\'i:
\href{https://github.com/BlackArch/blackarch}{https://github.com/BlackArch/blackarch}.
Tambi\'en hay varios repositorios secundarios aqu\'i:
\href{https://github.com/BlackArch}{https://github.com/BlackArch}.

En el repositorio git principal, hay tres directorios importantes:

\begin{itemize}
\item docs - Documentaci\'on.
\item packages - archivos PKGBUILD.
\item scripts - Scripts poco \'utiles.
\end{itemize}

\subsection{Scripts}
Aqu\'i hay una referencia para los scripts en el directorio \verb|scripts/| :

\begin{itemize}
\item baaur -  Pronto, esto subir\'a paquetes al AUR.
\item babuild - Construye un paquete.
\item bachroot - Administrar un chroot para pruebas.
\item baclean - Limpie los archivos.pkg.tar.xz antiguos del paquete repo.
\item baconflict - Pronto esto reemplazar\'a a \verb|scripts/conflicts|.
\item bad-files - Encuentra archivos defectuosos en los paquetes construidos.
\item balock - Obtenga o libere el bloqueo de reposici\'on de paquetes.
\item banotify - Notificar a IRC sobre el envio de paquetes.
\item barelease - Liberar paquetes al repositorio de paquetes.
\item baright - Imprimir la informaci\'on de copyright de BlackArch.
\item basign -  Firmar paquetes.
\item basign-key - Firmar una clave.
\item blackman - Esto se comporta como pacman pero construye a partir de git (no para de ser
    confundido con el Blackman de nrz).
\item check-groups - Verifica grupos.
\item checkpkgs - Verifica paquetes con errores.
\item conflicts - Compruebe si hay conflictos de archivos.
\item dbmod - Modificar una base de datos de paquetes.
\item depth-list - Cree una lista ordenada por la profundidad de la relaci\'on.
\item deptree - Cree un \'arbol de dependencias con una lista de solo los paquetes proporcionados por Blackarch.
\item get-blackarch-deps - Obtenga una lista de dependencias de Blackarch para un paquete.
\item get-official - Obtenga paquetes oficiales para su publicaci\'on.
\item list-loose-packages - Lista de paquetes que no est\'an en grupos y que no est\'an
    para otros paquetes.
\item list-needed - Enumerar las dependencias que faltan.
\item list-removed - Lista los paquetes que est\'an en el repositorio de paquetes pero no en git.
\item list-tools - Lista de herramientas.
\item outdated - Busque paquetes que est\'en desactualizados en el repositorio de paquetes.
    en relaci\'on con el git repo.
\item pkgmod - Modificar un paquete de compilaci\'on.
\item pkgrel - Incrementar pkgrel en un paquete.
\item prep - Limpia el estilo de un archivo PKGBUILD y encuentra errores.
\item sitesync - Sincronizaci\'on entre una copia local del repositorio de paquetes y una copia remota.
\item size-hunt - B\'usqueda de paquetes grandes.
\item source-backup - Haga copias de seguridad de los archivos de origen de los paquetes.
\end{itemize}

\section{Contributing to repository}
Esta secci\'on muestra c\'omo contribuir al proyecto BlackArch Linux. Nosotros aceptamos pull request de todos los tamaños, desde pequeñas correcciones de errores tipogr\'aficos hasta nuevos paquetes.
\\Para ayudar, sugerencias o preguntas no dude en ponerse en contacto con nosotros.
\\\\
Todos son Bienvenidos a contribuir. Todas las contribuciones son apreciadas.

\subsection{Tutoriales necesarios}
Por favor lea los siguientes tutorial antes de empezar a colaborar:
\begin{itemize}
\item
\href{https://wiki.archlinux.org/index.php/Arch\_Packaging\_Standards)}{Normas
de empaquetado de Arch}
\item \href{https://wiki.archlinux.org/index.php/Creating\_Packages}{Creaci\'on 
de Paquetes}
\item \href{https://wiki.archlinux.org/index.php/PKGBUILD}{PKGBUILD}
\item \href{https://wiki.archlinux.org/index.php/Makepkg}{Makepkg}
\end{itemize}

\subsection{Steps for contributing}
Para enviar sus cambios al proyecto BlackArchLinux, siga estos pasos
pasos:
\begin{enumerate}
\item Fork al repositorio origen
\url{https://github.com/BlackArch/blackarch}
\item Modifique los archivos necesarios (por ejemplo, PKGBUILD, archivos.patch, etc.).
\item Commit a tus cambios.
\item Sube tus cambios.
\item P\'idanos que fusionemos sus cambios, preferiblemente a trav\'es de una solicitud de pull request
\end{enumerate}

\subsection{Ejemplo}
El siguiente ejemplo muestra c\'omo enviar un nuevo paquete a BlackArch
proyecto. Utilizamos \href{https://wiki.archlinux.org/index.php/yaourt}{yaourt}
(tambi\'en puede usar pacaur) para obtener un archivo PKGBUILD preexistente 
\textbf{nfsshell} de \href{https://aur.archlinux.org/}{AUR} y ajustarlo de acuerdo a nuestras necesidades.

\subsubsection{Sincronizar PKGBUILD}
Sincronizar el fichero \textit{PKGBUILD} usando yaourt o pacaur:
\begin{lstlisting}
  user@blackarchlinux $ yaourt -G nfsshell
  ==> Download nfsshell sources
  x LICENSE
  x PKGBUILD
  x gcc.patch
  user@blackarchlinux $ cd nfsshell/
\end{lstlisting}

\subsubsection{Limpiar PKGBUILD}
Limpia el archivo \textit{PKGBUILD} y ahorra tiempo:
\begin{lstlisting}
  user@blackarchlinux nfsshell $ ./blackarch/scripts/prep PKGBUILD
  cleaning 'PKGBUILD'...
  expanding tabs...
  removing vim modeline...
  removing id comment...
  removing contributor and maintainer comments...
  squeezing extra blank lines...
  removing '|| return'...
  removing leading blank line...
  removing $pkgname...
  removing trailing whitespace...
\end{lstlisting}

\subsubsection{Ajustar PKGBUILD}
Ajustar el archivo \textit{PKGBUILD} :
\begin{lstlisting}
  user@blackarchlinux nfsshell $ vi PKGBUILD
\end{lstlisting}

\subsubsection{Construcci\'on del paquete}
Construcci\'on del paquete:
\begin{lstlisting}user@blackarchlinux nfsshell $ makepkg -sf
==> Making package: nfsshell 19980519-1 (Mon Dec  2 17:23:51 CET 2013)
==> Checking runtime dependencies...
==> Checking buildtime dependencies...
==> Retrieving sources...
-> Downloading nfsshell.tar.gz...
% Total    % Received % Xferd  Average Speed   Time    Time     Time
CurrentDload  Upload   Total   Spent    Left  Speed100 29213  100 29213    0
0  48150      0 --:--:-- --:--:-- --:--:-- 48206
-> Found gcc.patch
-> Found LICENSE
...
<lots of build process and compiler output here>
...
==> Leaving fakeroot environment.
==> Finished making: nfsshell 19980519-1 (Mon Dec  2 17:23:53 CET 2013)
\end{lstlisting}

\subsubsection{Instalar y probar el paquete}
Instalar y probar el paquete:
\begin{lstlisting}
  user@blackarchlinux nfsshell $ pacman -U nfsshell-19980519-1-x86_64.pkg.tar.xz
  user@blackarchlinux nfsshell $ nfsshell # test it
\end{lstlisting}

\subsubsection{dAd, commit and push package}
Add, commit y push el paquete
\begin{lstlisting}user@blackarchlinux nfsshell $ cd /blackarchlinux/packages
user@blackarchlinux ~/blackarchlinux/packages $ mv ~/nfsshell .
user@blackarchlinux ~/blackarchlinux/packages $ git commit -am nfsshell && git push
\end{lstlisting}

\subsubsection{Crea un pull request}
Crea una pull request en \href{https://github.com/}{github.com}
\begin{lstlisting}
  firefox https://github.com/<contributor>/blackarchlinux
\end{lstlisting}

\subsubsection{Adding a remote for upstream}
Una cosa inteligente a hacer si usted est\'a trabajando upstream y en un fork es tirar de su propio fork y añadir el repositorio principal como remoto
\begin{lstlisting}
  user@blackarchlinux ~/blackarchlinux $ git remote -v
  origin <the url of your fork> (fetch)
  origin <the url of your fork> (push)
  user@blackarchlinux ~/blackarchlinux $ git remote add upstream https://github.com/blackarch/blackarch
  user@blackarchlinux ~/blackarchlinux $ git remote -v
  origin <the url of your fork> (fetch)
  origin <the url of your fork> (push)
  upstream https://github.com/blackarch/blackarch (fetch)
  upstream https://github.com/blackarch/blackarch (push)
\end{lstlisting}

Por defecto, git deber\'ia empujar directamente al origen, pero aseg\'urate de que tu configuraci\'on de git est\'a correctamente. Esto no ser\'a un problema a menos que tengas derechos como
no ser\'as capaz de empujar upstream sin ellos.

Si tienes la capacidad de comprometerte, puedes tener m\'as \'exito usando
\textit{git@github.com:blackarch/blackarch.git} pero depende de ti.

\subsection{Solicitudes}
\begin{enumerate}
\item \textbf{Responsable} o \textbf{Colaborador} no añada comentarios en los
archivos \textit{PKGBUILD}. Añada sus nombres de responsable y colaborador
en la seccion de AUTORES de la BlackArch guide.
\item En favor de la coherencia, por favor, siga el estilo general de la otra
\textit{PKGBUILD} en el repositorio y usar la sangr\'ia de dos espacios.
\end{enumerate}


\subsection{Consejos generales}
\href{http://wiki.archlinux.org/index.php/Namcap}{namcap} puede comprobar si 
los paquetes contienen errores.

%------------------%
%  Capitulo 4      %
%------------------%

\chapter{Gu\'ia de Herramientas}
Pr\'oximamente...

\section{Pr\'oximamente...}
Pr\'oximamente...

%%% Apendice %%%
\appendix
%%%%%%%%%%%%%%%%%%%%%%%%%%%%%%%%%%%%%%%%%%%%%%%%%%%%%%%%%%%%%%%%%%%%%%%%%%%%%%%%
%                                                                              %
% BlackArch Linux plantilla de ap\'endice                                      %
%                                                                              %
%%%%%%%%%%%%%%%%%%%%%%%%%%%%%%%%%%%%%%%%%%%%%%%%%%%%%%%%%%%%%%%%%%%%%%%%%%%%%%%%

\appendix

\chapter{Ap\'endice}

\section{FAQ's}

\section{AUTORES}
Las siguientes personas han colaborado directamente en BlackArch:
\begin{itemize}
\item Tyler Bennnett (tylerb@trix2voip.com)
\item fnord0 (fnord0@riseup.net)
\item nrz (nrz@nullsecurity.net)
\item Ellis Kenyo (elken.tdos@gmail.com)
\item CaledoniaProject (the.warl0ck.1989@gmail.com)
\item sudokode (sudokode@gmail.com)
\item Valentin Churavy (v.churavy@gmail.com)
\item Boy Sandy Gladies Arriezona (reno.esper@gmail.com)
\item Mathias Nyman
\item Johannes Löthberg (demizide@gmail.com)
\item Thiago da Silva Teixeira (teixeira.zeus@gmail.com)
\end{itemize}
Las siguientes personas han colaborado directamente en ArchPwn, 
y se han unido a BlackArch:
\begin{itemize}
\item Francesco Piccinno (stack.box@gmail.com)
\item jensp (jens@jenux.homelinux.org)
\item Valentin Churavy (v.churavy@gmail.com)
\end{itemize}
Hemos tomado build code de las siguientes personas:
\begin{itemize}
\item  3ED (krzysztof1987@gmail.com)
\item  AUR Perl (aurperl@juster.info)
\item  Aaron Griffin (aaron@archlinux.org)
\item  Abakus (java5@arcor.de)
\item  Adam Wolk (netprobe@gmail.com)
\item  Aleix Pol (aleixpol@kde.org)
\item  Aleshus (aleshusi@gmail.com)
\item  Alessandro Pazzaglia (jackdroido@gmail.com)
\item  Alessandro Sagratini (ale\_sagra@hotmail.com)
\item  Alex Cartwright (alexc223@googlemail.com)
\item  Alexander De Sousa (archaur.xandy21@spamgourmet.com)
\item  Alexander Rødseth (rodseth@gmail.com)
\item  Allan McRae (allan@archlinux.org)
\item  AmaN (gabroo.punjab.da@gmail.com)
\item  Andre Klitzing (aklitzing@online.de)
\item  Andrea Scarpino (andrea@archlinux.org)
\item  Andreas Schönfelder (passtschu@freenet.de)
\item  Andrej Gelenberg (andrej.gelenberg@udo.edu)
\item  Angel Velasquez (angvp@archlinux.org)
\item  Antoine Lubineau (antoine@lubignon.info)
\item  Anton Bazhenov (anton.bazhenov@gmail.com)
\item  Arkham (arkham@archlinux.us)
\item  Arthur Danskin (arthurdanskin@gmail.com)
\item  Balda (balda@balda.ch)
\item  Balló György (ballogyor+arch@gmail.com)
\item  Bartek Piotrowski (barthalion@gmail.com)
\item  Bartosz Feński (fenio@debian.org)
\item  Bartłomiej Piotrowski (nospam@bpiotrowski.pl)
\item  Bogdan Szczurek (thebodzio@gmail.com)
\item  Brad Fanella (bradfanella@archlinux.us)
\item  Brian Bidulock (bidulock@openss7.org)
\item  C Anthony Risinger (anthony@xtfx.me)
\item  CRT (crt.011@gmail.com)
\item  Can Celasun (dcelasun@gmail.com)
\item  Chaniyth (chaniyth@yahoo.com)
\item  Chris Brannon (cmbrannon79@gmail.com)
\item  Chris Giles (Chris.G.27@gmail.com) \& daschu117
\item  Christoph Siegenthaler (csi@gmx.ch)
\item  Christoph Zeiler (archNOSPAM@moonblade.org)
\item  Clément DEMOULINS (clement@archivel.fr)
\item  Corrado Primier (bardo@aur.archlinux.org)
\item  Daenyth (Daenyth+Arch@gmail.com)
\item  Dale Blount (dale@archlinux.org)
\item  Damir Perisa (damir.perisa@bluewin.ch)
\item  Dan Fuhry (dan@fuhry.us)
\item  Dan Serban (dserban01@yahoo.com)
\item  Daniel A. Campoverde Carrión
\item  Daniel Golle
\item  Daniel Griffiths (ghost1227@archlinux.us)
\item  Daniel J Griffiths (ghost1227@archlinux.us)
\item  Daniel Micay (danielmicay@gmail.com)
\item  Dave Reisner (dreisner@archlinux.org)
\item  Dawid Wrobel (cromo@klej.net)
\item  Devaev Maxim (mdevaev@gmail.com)
\item  Devin Cofer (ranguvar@archlinux.us)
\item  DigitalPathogen (aur@InfoSecResearchLabs.co.uk)
\item  DigitalPathogen (aur@digitalpathogen.co.uk)
\item  Dmitry A. Ilyashevich (dmitry.ilyashevich@gmail.com)
\item  Dominik Heidler (dheidler@gmail.com)
\item  DrZaius (lou@fakeoutdoorsman.com)
\item  Ebubekir KARUL (ebubekirkarul@yandex.com)
\item  Eduard "bekks" Warkentin (eduard.warkentin@gmail.com)
\item  Elmo Todurov (todurov@gmail.com)
\item  Emmanuel Gil Peyrot (linkmauve@linkmauve.fr)
\item  Eric Belanger (eric@archlinux.org)
\item  Ermak (ermak@email.it)
\item  Evangelos Foutras (evangelos@foutrelis.com)
\item  Fabian Melters (melters@gmail.com)
\item  Fabiano Furtado (fusca14@gmail.com)
\item  Federico Quagliata (ntp@quaqo.org)
\item  Firmicus (francois.archlinux@org)
\item  Florian Pritz (bluewind@jabber.ccc.de)
\item  Florian Pritz (flo@xinu.at)
\item  Francesco Piccinno (stack.box@gmail.com)
\item  François Charette (francois@archlinux.org)
\item  Gaetan Bisson (bisson@archlinux.org)
\item  Geoffroy Carrier (geoffroy.carrier@koon.fr)
\item  Georg Grabler (STiAT)
\item  George Hilliard (gh403@msstate.edu)
\item  Gerardo Exequiel Pozzi (vmlinuz386@yahoo.com.ar)
\item  Gilles CHAUVIN (gcnweb@gmail.com)
\item  Giovanni Scafora (giovanni@archlinux.org)
\item  Gordin (9ordin@gmail.com)
\item  Guillaume ALAUX (guillaume@archlinux.org)
\item  Guillermo Ramos (0xwille@gmail.com)
\item  Gustavo Alvarez (sl1pkn07@gmail.com)
\item  Hugo Doria (hugo@archlinux.org)
\item  Hyacinthe Cartiaux (hyacinthe.cartiaux@free.fr)
\item  James Fryman (jfryman@gmail.com)
\item  Jan "heftig" Steffens (jan.steffens@gmail.com)
\item  Jan de Groot (jgc@archlinux.org)
\item  Jaroslav Lichtblau (dragonlord@aur.archlinux.org)
\item  Jaroslaw Swierczynski (swiergot@aur.archlinux.org)
\item  Jason Chu (jason@archlinux.org)
\item  Jason R Begley (jayray@digitalgoat.com)
\item  Jason Rodriguez
\item  Jason St. John (jstjohn@purdue.edu)
\item  Jawmare (victor2008@gmail.com)
\item  Jeff Mickey (jeff@archlinux.org)
\item  Jens Pranaitis (jens@chaox.net)
\item  Jens Pranaitis (jens@jenux.homelinux.org)
\item  Jinx (jinxware@gmail.com)
\item  John D Jones III (jnbek1972@gmail.com)
\item  John Proctor (jproctor@prium.net)
\item  Jon Bergli Heier (snakebite@jvnv.net)
\item  Jonas Heinrich
\item  Jonathan Steel (jsteel@aur.archlinux.org)
\item  Joris Steyn (jorissteyn@gmail.com)
\item  Josh VanderLinden (arch@cloudlery.com)
\item  Jozef Riha (jose1711@gmail.com)
\item  Judd Vinet (jvinet@zeroflux.org)
\item  Juergen Hoetzel (jason@archlinux.org)
\item  Juergen Hoetzel (juergen@archlinux.org)
\item  Justin Davis (jrcd83@gmail.com)
\item  Kaiting Chen (kaitocracy@gmail.com)
\item  Kaos
\item  Kevin Piche (kevin@archlinux.org)
\item  Kory Woods (kory@virlo.net)
\item  Kyle Keen (keenerd@gmail.com)
\item  Larry Hajali (larryhaja@gmail.com)
\item  LeCrayonVert
\item  Le\_suisse (lesuisse.dev+aur@gmail.com)
\item  Lekensteyn (lekensteyn@gmail.com)
\item  Limao Luo (luolimao+AUR@gmail.com)
\item  Lucien Immink
\item  Lukas Fleischer (archlinux@cryptocrack.de)
\item  Manolis Tzanidakis
\item  Marcin "avalan" Falkiewicz (avalatron@gmail.com)
\item  Mariano Verdu (verdumariano@gmail.com)
\item  Marti Raudsepp (marti@juffo.org)
\item  MatToufoutu (mattoufootu@gmail.com)
\item  Matthew Sharpe (matt.sharpe@gmail.com)
\item  Mauro Andreolini (mauro.andreolini@unimore.it)
\item  Max Pray a.k.a. Synthead (synthead@gmail.com)
\item  Max Roder (maxroder@web.de)
\item  Maxwell Pray a.k.a. Synthead (synthead@gmail.com)
\item  Maxwell Pray a.k.a. Synthead (synthead1@gmail.com)
\item  Mech (tiago.bmp@gmail.com)
\item  Michael Düll (mail@akurei.me)
\item  Michael P (ptchinster@archlinux.us)
\item  Michal Krenek (mikos@sg1.cz)
\item  Michal Zalewski (lcamtuf@coredump.cx)
\item  Miguel Paolino (mpaolino@gmail.com)
\item  Miguel Revilla (yo@miguelrevilla.com)
\item  Mike Roberts (noodlesgc@gmail.com)
\item  Mike Sampson (mike@sambodata.com)
\item  Nassim Kacha (nassim.kacha@gmail.com)
\item  Nicolas Pouillard (nicolas.pouillard@gmail.com)
\item  Nicolas Pouillard https://nicolaspouillard.fr
\item  Niklas Schmuecker
\item  Oleander Reis (oleander@oleander.cc)
\item  Olivier Le Moal (mail@olivierlemoal.fr)
\item  Olivier Médoc "oliv" (o\_medoc@yahoo.fr)
\item  Pascal E. (archlinux@hardfalcon.net)
\item  Patrick Leslie Polzer (leslie.polzer@gmx.net)
\item  Paul Mattal (paul@archlinux.org)
\item  Paul Mattal (pjmattal@elys.com)
\item  Pengyu CHEN (cpy.prefers.you@gmail.com)
\item  Peter Wu (lekensteyn@gmail.com)
\item  Philipp 'TamCore' B. (philipp@tamcore.eu)
\item  Pierre Schmitz (pierre@archlinux.de)
\item  Pranay Kanwar (pranay.kanwar@gmail.com)
\item  Pranay Kanwar (warl0ck@metaeye.org)
\item  PyroPeter (googlemail@com.abi1789)
\item  PyroPeter (googlemail.com@abi1789)
\item  Ray Rashif (schiv@archlinux.org)
\item  Remi Gacogne
\item  Renan Fernandes (renan@kauamanga.com)
\item  Richard Murri (admin@richardmurri.com)
\item  Roberto Alsina (ralsina@kde.org)
\item  Robson Peixoto (robsonpeixoto@gmail.com)
\item  Roel Blaauwgeers (roel@ttys0.nl)
\item  Rorschach (r0rschach@lavabit.com)
\item  Ruben Schuller (shiml@orgizm.net)
\item  Rudy Matela (rudy@matela.com)
\item  Ryon Sherman (ryon.sherman@gmail.com)
\item  Sabart Otto \item  Seberm (seberm@gmail.com)
\item  SakalisC (chrissakalis@gmail.com)
\item  Sam Stuewe (halosghost@archlinux.info)
\item  SanskritFritz (SanskritFritz@gmail.com)
\item  Sarah Hay (sarahhay@mb.sympatico)
\item  Sebastian Benvenuti (sebastianbenvenuti@gmail.com)
\item  Sebastian Nowicki (sebnow@gmail.com)
\item  Sebastien Duquette (ekse.0x@gmail.com)
\item  Sebastien LEDUC (sebastien@sleduc.fr)
\item  Sebastien Leduc (sebastien@sleduc.fr)
\item  Sergej Pupykin (pupykin.s+arch@gmail.com)
\item  Sergio Rubio (rubiojr@biondofu.net)
\item  Sheng Yu (magicfish1990@gmail.com)
\item  Simon Busch (morphis@gravedo.de)
\item  Simon Legner (Simon.Legner@gmail.com)
\item  Sirat18 (aur@sirat18.de)
\item  SpepS (dreamspepser@yahoo.it)
\item  Spider.007 (archPackage@spider007.net)
\item  Stefan Seering
\item  Stephane Travostino (stephane.travostino@gmail.com)
\item  Stéphane Gaudreault (stephane@archlinux.org)
\item  Sven Kauber (celeon@gmail.com)
\item  Sven Schulz (omee@archlinux.de)
\item  Sébastien Duquette (ekse.0x@gmail.com)
\item  Sébastien Luttringer (seblu@archlinux.org)
\item  TDY (tdy@gmx.com)
\item  Teemu Rytilahti (tpr@iki.fi)
\item  Testuser\_01
\item  Thanx (thanxm@gmail.com)
\item  Thayer Williams (thayer@archlinux.org)
\item  Thomas S Hatch (thatch45@gmail.com)
\item  Thorsten Töpper
\item  Tilmann Becker (tilmann.becker@freenet.de)
\item  Timothy Redaelli (timothy.redaelli@gmail.com)
\item  Timothée Ravier (tim@siosm.fr)
\item  Tino Reichardt
\item  Tobias Kieslich (tobias@justdreams.de)
\item  Tobias Powalowski (tpowa@archlinux.org)
\item  Tom K (tomk@runbox.com)
\item  Tom Newsom (Jeepster@gmx.co.uk)
\item  Tomas Lindquist Olsen (tomas@famolsen.dk)
\item  Travis Willard (travisw@wmpub.ca)
\item  Valentin Churavy (v.churavy@gmail.com)
\item  ViNS (gladiator@fastwebnet.it)
\item  Vlatko Kosturjak (kost@linux.hr)
\item  Wes Brown (wesbrown18@gmail.com)
\item  William Rea (sillywilly@gmail.com)
\item  Xavier Devlamynck (magicrhesus@ouranos.be)
\item  Xiao\item Long Chen (chenxiaolong@cxl.epac.to)
\item  aeolist (aeolist@hotmail.com)
\item  ality@pbrane.org
\item  astaroth (astaroth\_@web.de)
\item  bender02@archlinux.us
\item  billycongo (billycongo@gmail.com)
\item  bslackr (brendan@vastactive.com)
\item  cbreaker (cbreaker@tlen.pl)
\item  chimeracoder (dev@chimeracoder.net)
\item  damir (damir@archlinux.org)
\item  danitool
\item  darkapex (me@jailuthra.in)
\item  daronin
\item  dkaylor (dpkaylor@gmail.com)
\item  dobo (dobo90\_at\_gmail@com)
\item  dorphell (dorphell@archlinux.org)
\item  evr (evanroman@at.gmail)
\item  fnord0 (fnord0@riseup.net)
\item  fxbru (frxbru@gmail)
\item  hcar
\item  icarus (icarus.roaming@gmail.com)
\item  iceman (icemanf@gmail.com)
\item  kastor (kastor@fobos.org)
\item  kfgz (kfgz@interia.pl)
\item  linuxSEAT (linuxSEAT@gmail.com)
\item  m4xm4n (max@maxfierke.com)
\item  mar77i (mysatyre@gmail.com)
\item  marc0s (marc0s@fsfe.org)
\item  mickael9 (mickael9@gmail.com)
\item  nblock (nblock@archlinux.us)
\item  nofxx (x@nofxx.com)
\item  onny (onny@project
\item  pootzko (pootzko@gmail.com)
\item  revel (revel@muub.net)
\item  rich\_o (rich\_o@lavabit.com)
\item  s1gma (s1gma@mindslicer.com)
\item  sandman (r.coded@gmail.com)
\item  sebikul (sebikul@gmail.com)
\item  sh0 (mee@sh0.org)
\item  shild (sxp@bk.ru)
\item  simo (simo@archlinux.org)
\item  snuo
\item  sudokode (sudokode@gmail.com)
\item  tobias (tobias@archlinux.org)
\item  trashstar (trash@ps3zone.org)
\item  unexist (unexist@subforge.org)
\item  untitled (rnd0x00@gmail.com)
\item  virtuemood (virtue@archlinux.us)
\item  wido (widomaker2k7@gmail.com)
\item  wodim (neikokz@gmail.com)
\item  yannsen (ynnsen@gmail.com)
\end{itemize}


\end{document}

%%% EOF %%%
